\section{Higher moments and algebraic consequences}\label{sec:moments}
This appendix extends the count used in the main construction to higher moments
and derives region and subdivision bounds. The probability argument uses classical
shattering and higher-gap methods; the geometric consequences use established
bounds for polynomial sign conditions and cylindrical algebraic decomposition.
They do not assert new general smoothing or real-algebraic algorithms.

\subsection{All fixed moments, including atomic noise}
Retain the notation $A_a$ and the interval-mass assumption
\eqref{eq:intervalmass} from Section~\ref{sec:enumeration}.

\begin{lemma}\label{lem:moments}
For $m\ge1$, let $D_a$ be the largest size of a coordinate set shattered
by $A_a$. For $1\le r\le m$ and every real $p\ge1$,
\begin{align}
 \Prob\{D_a\ge r\}&\le \binom{m}{r}(2\phi a+\tau)^r,
                                                        \label{eq:VCtail}\\
 \E|A_a|^p&\le[1+(m+1)^p(2\phi a+\tau)]^m
 \le\exp\{m(m+1)^p(2\phi a+\tau)\}.\label{eq:allmoments}
\end{align}
\end{lemma}
\begin{proof}
The conditional class-minimum argument in Lemma~\ref{lem:count} bounds
the probability that a particular $r$-set is shattered by
$(2\phi a+\tau)^r$. A larger shattered set contains a shattered $r$-set,
so a union bound proves \eqref{eq:VCtail}. The classical Sauer bound follows
here directly from \eqref{eq:sandwich}: every shattered set has size at
most $D_a$, whence
\[
 |A_a|\le\sum_{j=0}^{D_a}\binom{m}{j}\le(m+1)^{D_a}.
\]
For the last inequality, encode a subset by its increasing elements padded
with zeros to length $D_a$; for $D_a=0$ both sides are one.
Set $b=(m+1)^p$. The integer-valued tail identity gives
\[
 \E b^{D_a}=1+\sum_{r=1}^m(b^r-b^{r-1})\Prob(D_a\ge r)
 \le\sum_{r=0}^m\binom{m}{r}[b(2\phi a+\tau)]^r.
\]
The binomial formula and $1+x\le e^x$ prove the claim.
\end{proof}

These are moment bounds for one shared random linear penalty, with arbitrary
deterministic support costs. Polynomial moments in smoothed optimization
are established \cite{RoglinTeng,BrunschRoglin}; the proof above records the
particular interval-mass version needed here.

\begin{theorem}[Moments on a compact parameter domain]\label{thm:parammoments}
Let $q_z$ be deterministic $L$-Lipschitz functions on a nonempty compact
metric space $T$, and let $K$ count all labels minimizing
$F_z(t)=q_z(t)+\xi^Tz$ somewhere on $T$. Fix real $p\ge1$ and
$n\ge\max\{1,m\}$. If
\begin{equation}\label{eq:momentdelta}
 \tau\le\frac1{2n(n+1)^p},\qquad
 \delta_p=\frac1{4n\phi(n+1)^p},
\end{equation}
and $T$ has a deterministic $\delta_p/(2L)$-net of size $J_p$ for $L>0$,
then $\E K^p\le eJ_p^p$. For $L=0$ use one point. On $[-R,R]^k$,
with $k\ge1$ and $R>0$, one may take
\begin{equation}\label{eq:momentnet}
 J_p\le[1+8kRLn\phi(n+1)^p]^k.
\end{equation}
For continuous noise with densities at most $\phi$, one may instead use
$\delta_p=1/[2n\phi(n+1)^p]$. For the grid \eqref{eq:noisegrid},
\eqref{eq:momentdelta} holds when $N\ge2n(n+1)^p$.
\end{theorem}
\begin{proof}
The exponent in \eqref{eq:allmoments}, at $a=\delta_p$, is at most one.
Every active support belongs to $A_{\delta_p}(t_j)$ at some net point,
by the three-term Lipschitz comparison in Theorem~\ref{thm:dictionary}.
Convexity therefore gives
\[
 K^p\le\left(\sum_j|A_{\delta_p}(t_j)|\right)^p
 \le J_p^{p-1}\sum_j|A_{\delta_p}(t_j)|^p.
\]
Taking expectations proves the bound without independence between net
points. The midpoint construction \eqref{eq:netsize}, with $\delta_p$
in place of $\varepsilon$, proves \eqref{eq:momentnet}. For $L=0$ every
branch is constant on $T$; tied constants are still counted separately.
Singleton domains and $m=0$ are immediate. Compactness proves measurability
as before Theorem~\ref{thm:dictionary}. Setting $\tau=0$ permits the stated
larger tolerance. The grid assertion follows from \eqref{eq:gridmass}.
\end{proof}

If $K_1,\ldots,K_s$ are message counts with this common bound, then
\begin{equation}\label{eq:sum-moments}
 \E\left(\sum_{j=1}^s K_j\right)^p
 \le s^{p-1}\sum_{j=1}^s\E K_j^p\le e(sJ_p)^p.
\end{equation}
Overlapping subtrees and shared noise do not create an independence
assumption. Taking $N$ to be the smallest admissible power of two needs
$O_p(\log(n+1))$ random bits per coordinate. This moment order, and hence
the grid, must be fixed before drawing the perturbation.

\subsection{A companion affine-rank gap estimate}
The following alternative tail estimate uses the classical higher-gap
conditioning method of R\"oglin and Teng \cite[Sections~6.1--6.2]{RoglinTeng}.
It is useful independently of shattering, but is not needed by the main
algorithm. The affine rank of a set is the dimension of its affine hull.

\begin{proposition}\label{prop:rankgap}
For $1\le r\le m$, independent coordinates with densities at most $\phi$
satisfy
\begin{equation}\label{eq:rankgap}
 \Prob\{|A_a|\ge2^{r-1}+1\}
 \le \binom{m}{r}2^{r(r+1)}(2\phi a)^r.
\end{equation}
Under \eqref{eq:intervalmass}, the same bound holds with the last factor
replaced by $(2r!\phi a+\tau)^r$.
\end{proposition}
\begin{proof}
A $d$-dimensional affine subspace admits an injective projection onto some
$d$ coordinates, so it contains at most $2^d$ cube points. The event in
\eqref{eq:rankgap} therefore supplies $r+1$ affinely independent supports
and $r$ coordinates $I$ on which their differences have full rank.
Their patterns $y_0,\ldots,y_r\in\{0,1\}^r$ form at most
$2^{r(r+1)}$ ordered lists. Condition outside $I$ and define the class
costs $C_y$ as in Lemma~\ref{lem:count}. Each selected class minimum is
in $[v,v+a]$, so
\[
 |(y_j-y_0)^T\xi_I+C_{y_j}-C_{y_0}|\le a
 \quad(1\le j\le r).
\]
The difference matrix $D$ is a nonsingular integer matrix with entries in
$\{-1,0,1\}$ and $|\det D|\ge1$. This event lies in a parallelepiped
of volume at most $(2a)^r$, giving the density bound. For atomic noise,
each cofactor of $D$ has absolute value at most $(r-1)!$, by its determinant
expansion. Every row sum of $|D^{-1}|$ is at most $r!$. The parallelepiped
is thus contained in an axis-parallel box of side lengths $2r!a$.
Conditional independence and \eqref{eq:intervalmass} give its probability
bound. Union over $I$ and the pattern lists finishes both proofs.
\end{proof}

\subsection{Connected regions for quadratic branches}
For a quadratic family on $T=[-R,R]^k$, $k\ge1$, $R>0$, let
$\mathcal A=\bigcup_{t\in T}\argmin_zF_z(t)$ and $K=|\mathcal A|$.
Let $\mathcal R$ count the connected components of the open sets
\[
 U_z=\{t\in(-R,R)^k:F_z(t)<F_w(t)\text{ for all }w\ne z\}.
\]
Thus one label can contribute several regions; a tie-only label contributes
none but remains in $\mathcal A$.

\begin{corollary}\label{cor:atomicregions}
For each fixed $k$ there is a constant $C_k$ such that, deterministically,
$\mathcal R\le C_kK^{k+1}$. Under Theorem~\ref{thm:parammoments} with
$p=k+1$, therefore,
\begin{equation}\label{eq:atomicregions}
 \E\mathcal R\le C_keJ_{k+1}^{k+1}.
\end{equation}
In particular, the endpoint grid with
$N\ge2n(n+1)^{k+1}$ gives polynomial expected connected-region count
in every fixed dimension, including $p=3$ in the plane.
\end{corollary}
\begin{proof}
Deleting never-active labels preserves the complete argmin set at every
point, not just its value. For $z\in\mathcal A$, the set $U_z$ is one
strict sign condition of the $K-1$ quadratic differences $F_w-F_z$
and the $2k$ linear box inequalities. The classical bound of Basu, Pollack,
and Roy \cite[Theorem~5.1]{BasuSigns} bounds its number of components by
$C_k(K+2k)^k$. Their theorem bounds the total number of components of all
realizable sign conditions, so it also bounds this one condition; a ball
containing the box supplies their bounded ambient convention. No general
position is assumed. An identically zero difference simply makes this
strict sign condition empty. Summing over $z$, and absorbing $2k$ into
a dimension-dependent constant since $K\ge1$, proves the deterministic
bound. Apply Theorem~\ref{thm:parammoments} for its expectation.
\end{proof}

Tie-only labels are essential to this reduction. The two planar branches
$0$ and $t_1^2$ have two open unique-support regions separated by their tie
line. Deleting the touching branch leaves one region. More generally, a
bound on labels alone cannot control regions for arbitrary Lipschitz
functions, which can cross infinitely often. Polynomial structure makes
\eqref{eq:atomicregions} possible. This coarse algebraic bound permits atomic
noise but does not preserve the sharper continuous constants proved in
Appendix~\ref{sec:geometry}.

\subsection{Constructing a subdivision with every tied label}
A sign-invariant cylindrical algebraic decomposition (CAD) partitions real
space into cells on each of which specified polynomials have constant signs.
Its exact construction has polynomial bit complexity for fixed dimension
and degree \cite[Introduction]{ArnonCAD}. We use that standard algorithm
without requiring adjacency information or a construction of maximal
connected regions.

\begin{corollary}\label{cor:subdivision}
Under the spectral assumptions of Theorem~\ref{thm:main}, fix the width.
There is a finite-grid algorithm that constructs, for every requested message,
a CAD subdivision of its closed box and lists \emph{all} minimizing support
labels on every cell. It is exact for every draw and has polynomial expected
bit complexity in the same numerical parameters as the main theorem.
It uses $O_p(\log(n+1))$ random bits per coordinate, with a fixed moment
order $p$ depending on the chosen CAD algorithm and the width.
\end{corollary}
\begin{proof}
Choose the integer $p\ge1$ below before sampling, and take $N$ to be the
smallest power of two satisfying $N\ge2n(n+1)^p$ in
\eqref{eq:noisegrid}. Rerun
Theorem~\ref{thm:dictionary} at accuracy and enumeration tolerance
\[
 \varepsilon_p=\delta=\frac1{12n\phi(n+1)^p}.
\]
On the $\varepsilon_p/(2L)$-net the emitted set satisfies
$E_j\subseteq A_{3\varepsilon_p}(t_j)$, and
\eqref{eq:allmoments} gives $\E|E_j|^p\le e$. The union $\mathcal D$
contains every active label and consists entirely of genuine support
branches, with
\begin{equation}\label{eq:enum-moments}
 \E|\mathcal D|^p\le e\widehat J_p^p,\qquad
 \widehat J_p\le[1+24kRLn\phi(n+1)^p]^k.
\end{equation}
Use one point in the constant or zero-dimensional cases.

Construct a sign-invariant CAD for every pairwise difference of dictionary
quadratics and every box-boundary polynomial, omitting identically zero
polynomials while retaining all their support labels. Keep all cells in
the closed box, including cells of smaller dimension. At each cell's exact
sample point, identify the minimum and list every label attaining it.
The signs certify that this list is constant throughout that cell.
The complete-label property of $\mathcal D$ proves equality with the argmin
set of the original family, even on boundary and tie cells. Extra inactive
branches may refine the subdivision, but cannot alter its labels.

Uniform Schur-complement bit bounds were proved in Theorem~\ref{thm:main}.
For fixed $k$, standard exact CAD and the additional output-label work
therefore cost at most $P_k(\ell)(1+|\mathcal D|)^{a(k)}$, for a fixed
computable exponent $a(k)$. Choose an integer $p\ge\max\{1,a(k)\}$,
uniformly over separator dimensions. Convexity and \eqref{eq:enum-moments}
bound the expectation of this work. The restricted grid oracle at the new
accuracy still has deterministic polynomial cost in the stated numerical
parameters; enumeration has polynomial expected work as before. Sum over
the requested messages. Shared noise does not require independent dictionaries.
\end{proof}

The cells in this output need not be maximal connected winning regions.
The grid is larger than the $N\ge2n$ grid sufficient for the main dictionary
algorithm. Thus the subdivision conclusion uses the higher-moment theorem,
not just a first-moment size bound followed by a polynomial-time algorithm.
